\begin{proposition}[Vertical and logarithmic tangent bundles]\hypertarget{tag-01SN}{}\label{inv:D1.2}
\StatementRecord{inv:D1.2}{proved}{01SN}{\SeriesRef{inv:D1.1}}

There is a canonical representable proper flat holomorphic map
$\bar f:X\to\mathcal B$ with $f=\rho\bar f$. If
\[
 \mathcal T=T_X(-\log W),\qquad
 \mathcal V=\ker\bigl(\mathcal T\to
                   \bar f^*T_{\mathcal B}(-p_0)\bigr),
\]
then, for rank-two vector bundles $E,G$ on $\mathcal B$,
\begin{equation}
 \mathcal V=\bar f^*E,\quad E=\OO\oplus\OO(-1),\qquad
 G=R^1\bar f_*\OO_X=\OO\oplus\OO(-1),
\label{pub:I:eq:D-2}
\end{equation}
\begin{equation}
 \hypertarget{tag-01OZ}{}\label{fam:eq:ab-direct-images}
R^0\bar f_*\OO_X=\OO,\qquad R^2\bar f_*\OO_X=\OO(-1).
\end{equation}
and the following is an exact sequence of locally free sheaves on $X$:
\begin{equation}
 \hypertarget{tag-01P0}{}\label{fam:eq:ab-log-tangent}
0\longrightarrow\bar f^*E\longrightarrow\mathcal T
 \longrightarrow\bar f^*\OO(-5)\longrightarrow0.       
\end{equation}
The trivial summands in $E,G$ are generated by the vertical field $\xi$
and the constant integral cocycle $\gamma$, respectively.
\end{proposition}



\begin{proof}
These bundles and maps are the standard-data calculation of
\cite[Lemma~6.2, pp.~74--75; Lemma~6.7, p.~78; Proposition~6.6, p.~77; Proposition~6.10, pp.~79--80]{DuEtAl2026v1}.
\begin{enumerate}[label=\arabic*.,leftmargin=*]
\item The fibre-divisor root objects and their proper flat
atlas comparison are those of \cref{push:src:E.3.1} with $k=l=1$.
\item On a smooth torus chart a vertical field is a holomorphic linear
combination of the translation fields. Its coefficients are constant
on each compact fibre and therefore depend only on the base. The
vertical derivative of an affine transition is its linear part $R_g$;
the translating part, including the logarithmic regluing, has vertical
derivative the identity. At the cusp the two fields
$x_1\partial_{x_1},x_2\partial_{x_2}$ extend as logarithmic vector
fields in every toric chart. The deck maps multiply each $x_i$ by a
base-dependent unit and a power of the base parameter. Their vertical
logarithmic derivative fixes both coefficient directions. The
exponential overlap sends $\partial_{\zeta_i}$ to
$2\pi i x_i\partial_{x_i}$. These transitions depend only on the base
and define $E$; they prove the bundle identity
$\mathcal V=\bar f^*E$.

\item At
$t=xy$, the kernel of $df$ is freely generated by
$x\partial_x-y\partial_y$ and $\partial_z$. At $t=xyz$ it is generated
by $x\partial_x-y\partial_y$ and $x\partial_x-z\partial_z$.
Indeed the equation $yz a+xz b+xy c=0$ in the regular local ring forces
$x\mid a,y\mid b,z\mid c$; after division the only relation is that
the three coefficients sum to zero. The logarithmic map sends
$\sum a_i z_i\partial_{z_i}$ to $(\sum a_i)t\partial_t$ and is
surjective. At a finite root point it is the ordinary relative tangent
sequence of the smooth map to the root coordinate $s$. These comparisons prove \eqref{fam:eq:ab-log-tangent}, since
$T_{\mathcal B}(-p_0)=\OO(7-12)$.

\item The invariant vector $e_2$ defines a subbundle $\OO\xi\subset E$.
It remains nonzero as a vector of the logarithmic vertical bundle at
the cusp, even where its image in the ordinary tangent fibre is zero.
The quotient line has automorphy factors $-1/\tau$ and $1/\tau$ under
the order-three and order-four generators. The modular discriminant
has weight twelve and no zero in the upper half-plane. Therefore its
reciprocal is a meromorphic section of the twelfth tensor power of
this quotient line. On the cusp chart
$\tau=(\log q)/(2\pi i)+h(q)$ with $h$ holomorphic, so the product
expansion of the discriminant is $q$ times a holomorphic unit.
There are no other zeros or poles, including on the finite root
charts. The quotient line thus has twelfth power
$\OO(-p_0)=\OO(-12)$, and the torsion-free Picard group in \ref{inv:D1.1} makes
it $\OO(-1)$. Its extension by $\OO$ splits because
$H^1(\mathcal B,\OO(1))=0$.

\item Proposition~\ref{push:src:E.3.1} gives local freeness,
geometric-fibre comparison and trace duality on the root atlas.
The determinant of \eqref{fam:eq:ab-log-tangent}, using $W=\bar f^*p_0$, gives
\begin{equation}
 \hypertarget{tag-01P2}{}\label{fam:eq:ab-relative-dualizing}
\omega_{X/\mathcal B}=\bar f^*(\det E)^{-1}.
\end{equation}
Relative Serre duality, applied on that same atlas and descended by
its trace pairing, gives $R^2\bar f_*\OO_X=\det E=\OO(-1)$.
Cup product is an isomorphism $\det G\to R^2\bar f_*\OO_X$, since it
is so on every geometric fibre by the preceding calculation.

\item The class $\gamma$ is nowhere zero in $G$. On a smooth marked fibre,
$H^1(\OO)$ is the quotient of constant lattice cocycles by the two
period rows. If $\gamma$ were a combination of those rows, evaluation
on $\widehat w,\widehat\delta$ would make both coefficients zero,
contradicting $\gamma(\widehat\gamma)=1$. The same argument applies
on each finite root cover, and affine translations act trivially on
these classes. At the cusp $\gamma$ kills
$\Lambda_{\rm tor}=\mathbb Z\langle\widehat w,\widehat\delta\rangle$.
Local sections of the free toric quotient with transitions
$s_i=\Psi_{\lambda_{ij}}s_j$ give the holomorphic constant \v Cech
cocycle $\gamma(\lambda_{ij})$ on the whole neighbourhood.
For $B_0=\left(\begin{smallmatrix}0&1\\-1&0\end{smallmatrix}\right)$,
the edges $e_1,e_2,e_2-e_1$ have deck vectors
$(0,1),(-1,0),(-1,-1)$, so its central values are $(0,-1,-1)$.
Their cup product with $(1,0,-1)$ is one by
\eqref{fam:eq:ab-cusp-resolution} and its displayed product.
On the punctured overlap the exponential map kills exactly
$\Lambda_{\rm tor}$, giving the same full-lattice cocycle.
The Leray edge map over the Stein cusp disc is an isomorphism,
since its positive coherent cohomology vanishes. This gives
a holomorphic nonvanishing extension of $[\gamma]$.
Thus
$\OO\gamma\subset G$ is a subbundle; its quotient is $\OO(-1)$ by
the determinant just computed. Once again $H^1(\OO(1))=0$ splits the
extension. This proves every assertion in \eqref{pub:I:eq:D-2}.
\end{enumerate}
\end{proof}